\documentclass[10pt,a4paper]{amsart}
\usepackage[margin=19mm]{geometry}
\usepackage{amsmath,amssymb,mathtools}
\usepackage[scr=rsfs,cal=euler]{mathalfa}
\usepackage{xfrac}
\usepackage[unicode,hidelinks,bookmarks=false]{hyperref}
\newcommand{\ie}{i.e.\ }
\newcommand{\cf}{cf.\ }
\newcommand{\resp}{respectively\ }
\newcommand{\Subsection}{\subsection}
\providecommand{\nobreakdash}{\nobreak\hbox{-}}
\setlength{\emergencystretch}{2em}
\multlinegap0pt
\usepackage{bm}
\usepackage{bbm}
\usepackage{dsfont}
\usepackage{xspace}
\usepackage{cancel}
\usepackage{mathtools}
\usepackage[shortlabels]{enumitem}
\usepackage{amsthm}
\makeatletter
\@ifundefined{theorem}{\newtheorem{theorem}{Theorem}}{}
\@ifundefined{lemma}{\newtheorem{lemma}[theorem]{Lemma}}{}
\makeatother
\DeclareMathOperator{\Syz}{Syz}
\newcommand{\KK}{\mathbb{K}}
\title[The complexity of the SupportMinors Modeling for the MinRank]{The complexity of the SupportMinors Modeling for the MinRank Problem}
\author{Daniel Cabarcas, Giulia Gaggero, Elisa Gorla}
\date{Source: arXiv:2506.06547v1 (CC BY 4.0)}
\begin{document}
\maketitle

\noindent\emph{Source and licence.} The complexity of the SupportMinors Modeling for the MinRank Problem. Version arXiv:2506.06547v1 (CC BY 4.0), distributed under the Creative Commons Attribution 4.0 International licence, as recorded in the arXiv licence metadata for this exact version and shown on the abstract page https://arxiv.org/abs/2506.06547v1. Full rights notice: licenses/arxiv-2506.06547-CC-BY-4.0.txt.\par\smallskip

\noindent\textit{Editorial scope.} Counted unit: the Theorem whose source label is \texttt{thm:supp} in the article (source0.tex lines 328--331, source label \texttt{thm:supp}), delivered together with the article's own complete original proof of it (source0.tex lines 333--367, one \texttt{proof} environment of 35 lines). No mathematics is written, repaired or re-derived: statement and proof are the article's own text line for line. Required context retained uncounted: lemma:supp (lines 286--293). Every definition, display and label of those lines is retained unchanged. Omitted from this package and not redistributed: 1--285, 294--327, 332--332, 368--526. Nothing inside a retained excerpt is omitted or altered: every retained body is delivered exactly as the article's lines are written, so no omission receipt of any kind accompanies this package. Citations: the retained excerpts carry no citation command at all, so no bibliography entry of the article is transcribed and no reference is redistributed. Reference adaptation: none was needed and none was made; every reference inside the delivered unit resolves inside it, so no unresolved reference or citation is produced. Printed slips kept literal: no printed word, symbol, hypothesis or number of the counted statement or of its proof was altered. Layout and class: the article's own class is \texttt{article} with packages including inputenc, amsfonts, amsmath, amsthm, amssymb, bm, wasysym, comment, enumerate, csquotes, babel, graphics, standalone, tikz; the wrapper is the lane's pinned \texttt{amsart} editorial wrapper at 10pt on A4, which restates the theorem-like environments the retained text needs and adds the article's own macro definitions that the retained text needs (\texttt{\textbackslash KK}, \texttt{\textbackslash Syz}), with extra wrapper packages (bm, bbm, dsfont, xspace, cancel, mathtools, enumitem). The delivered numbering is the unit's own. Assets: no retained span contains a figure environment, a graphics-inclusion command, a PDF-page inclusion, a table or a TikZ picture; no image file is copied into this package, the package declares no asset dependency and no image file is redistributed with it. Licence: version arXiv:2506.06547v1 of this article is distributed under the Creative Commons Attribution 4.0 International licence, as recorded in the arXiv licence metadata for this exact version and shown on the abstract page https://arxiv.org/abs/2506.06547v1. A full rights notice is delivered at licenses/arxiv-2506.06547-CC-BY-4.0.txt. Characters: every retained excerpt is pure ASCII, so the delivered engine, XeLaTeX with its default Latin Modern text font, reports no missing character.
\begin{lemma}\label{lemma:supp}
Each $S\in \mathbf{S}$ falls into one of these mutually exclusive cases:
\begin{enumerate}[i)]
\item $S$ is supported on the positions corresponding to $\mathcal{G}$.
\item $S$ is supported on the positions corresponding to $\mathcal{F}\setminus \mathcal{G}$.
\item $S$ does not fall into case i) or ii) and it involves a variable $c_{ij}$ in a position that corresponds to an element of $\mathcal{F}\setminus \mathcal{G}$.
\end{enumerate}  
\end{lemma}

\begin{theorem}\label{thm:supp}
Let $P=\KK[x_\ell,c_{ij}: \ell\in[K], i\in[r], j\in[n]]$ be an $R$-algebra with homomorphism $\rho:R\rightarrow P$ given by $c_{ij}\mapsto c_{ij}$ and $y_{kj}\mapsto m_{kj}$.
Let $T\in \Syz(\mathcal{F})$. If $\rho(T)\in \Syz(\rho(\mathcal{G}))\setminus\{0\}$, then $T\in\Syz(\mathcal{G})$.
\end{theorem}

\begin{proof} 
Let $\mathbf{T}_1$ be the set of elements of $\mathbf{S}$ that fall into case i) of Lemma~\ref{lemma:supp}, let $\mathbf{T}_2$ be the set of those that fall into case ii), and let $\mathbf{T}_3$ be the set of those that fall into case iii). By Lemma~\ref{lemma:supp}, $\mathbf{S}=\mathbf{T}_1\cup\mathbf{T}_2\cup\mathbf{T}_3$. 
Let $T\in \Syz(\mathcal{F})$ and suppose that $\rho(T)\in \Syz(\rho(\mathcal{G}))$.
Since $\mathbf{S}$ generates $\Syz(\mathcal{F})$, by Lemma~\ref{lemma:supp} we can write
\begin{equation}\label{eqn:expT}
T=\sum_{S\in\mathbf{T}_1}\alpha_S S + 
\sum_{S\in\mathbf{T}_2}\alpha_S S + 
\sum_{S\in\mathbf{T}_3} \alpha_S S.
\end{equation}
Up to replacing $T$ by $T-\sum_{S\in\mathbf{T}_1}\alpha_S S$, we may assume that $\alpha_S=0$ for every $S\in\mathbf{T}_1$. In particular,
\begin{equation}\label{eqn:rhoT}
\rho(T) = \sum_{S\in\mathbf{T}_2}\alpha_S \rho(S) + 
\sum_{S\in\mathbf{T}_3} \alpha_S\rho(S).
\end{equation}
Our thesis corresponds to proving that $T=0$. 

We start by discussing how the support changes when passing from $T$ to $\rho(T)$. Since $\rho(c_{ij})=c_{ij}$ for all $i$ and $j$, if one entry of $T$ involves a $c$-variable with a nonzero coefficient, then the corresponding entry of $\rho(T)$ involves the same $c$-variable with the same coefficient. In particular, when looking at the $c$-variables, the supports of $T$ and $\rho(T)$ coincide. Since $\rho(T)\in\Syz(\mathcal{G})$, the positions of $T$ corresponding to elements of $\mathcal{F}\setminus\mathcal{G}$ cannot involve any $c$ variable.

A coefficient $c_{ij}$ in position $(I,J)$ can only come from an element of $\mathbf{T}_2$ or $\mathbf{T}_3$ corresponding to the multisets $I\cup\{i\}$ and $J\cup\{j\}$. Therefore, if $c_{ij}E_{I,J}$ comes from a syzygy $\rho(S)$ and cancels in (\ref{eqn:expT}), then it cancels with a summand $c_{ij}E_{I,J}$ coming from a different syzygy $\rho(S^\prime)$ which corresponds to the same multisets $I\cup\{i\}$ and $J\cup\{j\}$. Therefore, we may restrict to $D_{I\cup\{i\},J\cup\{j\}}$ and only discuss which cancellations occur in (\ref{eqn:expT}) for syzygies that originate from it. 

Let $S\in\mathbf{T}_3$ and let $I,J$ be the multisets from which $S$ originates. Then $I\subseteq[m+r]$, $J\subseteq[n]$, $|I|=|J|=r+2$ and one of the following must hold:
\begin{enumerate}
\item $I$ is a set with $|I\cap[m]|=2$, $J$ contains one element $k$ with multiplicity two and every other element has multiplicity one,  and $S\in\mathbf{S}_2$.
\item $I$ and $J$ are sets of cardinality $r+2$, $|I\cap[m]|=2$, and $S\in\mathbf{S}_3$ with $h>2$.
\item $I$ and $J$ are sets of cardinality $r+2$, $|I\cap[m]|=2$, and $S\in\mathbf{S}_4$ with $k>1$.
\end{enumerate}

In case 1., there exists exactly one syzygy $\Sigma\in\mathbf{T}_3$ that originates from those $I$ and $J$. Write $I=\{i_1,\ldots,i_{r+2}\}$ with $i_1<\ldots<i_{r+2}$, $J=\{k,j_1,\ldots,j_{r+1}\}$ with $j_1<\ldots<j_{r+1}$ and $k\in\{j_1,\ldots,j_{r+1}\}$. In this case we have $\alpha_{\Sigma}=0$, since otherwise the element $\alpha_{\Sigma}c_{i_3 k}E_{I\setminus\{i_3\},\{j_1,\ldots,j_{r+1}\}}$ does not cancel in (\ref{eqn:expT}), as there is exactly one syzygy in $\mathbf{T}_2\cup \mathbf{T}_3$ where $c_{i_3 k}$ appears in position $I\setminus\{i_3\},\{j_1,\ldots,j_{r+1}\}$.

In cases 2. and 3., assume for ease of notation that $I=\{1,2,m+1,\ldots,m+r\}$ and $J=[r+2]$.
Because of what we discussed above, when studying cancellations among the $c$-variables, we may restrict our attention to the syzygies that originate from the $I$ and $J$ that we just fixed. For $h>2$, $c_{h-2,1}E_{I\setminus\{m+h-2\},\{2,\ldots,r+2\}}$ only appears in syzygies obtained from developing the determinant of $D_{I,J}$ with respect to row $h$ or column $1$. The only syzygy in $\mathbf{T}_2\cup \mathbf{T}_3$ which involves $c_{h-2,1}$ is $\Sigma_h=|D_{I,J}|_{1,\bullet} - |D_{I,J}|_{h,\bullet}$. Therefore no cancellation is possible and $\alpha_{\Sigma_h}=0$ for $h>2$. For a fixed $2\leq k\leq r+2$, $c_{1,k}E_{I\setminus\{m+1\},J\setminus\{k\}}$ only appears in syzygies obtained from developing the determinant of $D_{I,J}$ with respect to row $3$ or column $k$.
Since $\alpha_{\Sigma_3}=0$, the only syzygy in which it appears is $\Theta_k=|D_{I,J}|_{1,\bullet} - |D_{I,J}|_{\bullet,k}$. Again, no cancelation is possible, showing that $\alpha_{\Theta_k}=0$ for $k>1$.

This proves that, if $\rho(T)\in\Syz(\mathcal{G})$, then there is an expression of $T$ as in (\ref{eqn:expT}) that does not involve any element of $\mathbf{T}_3$. Since the support of $\rho(S)$ is disjoint from $\mathcal{G}$ for every $S\in\mathbf{T}_2$, we deduce that $\rho(T)\in\Syz(\mathcal{G})$ forces $T=0$. 
\end{proof}

\end{document}
